% ============================================
% Глава 4. Обсуждение и направления развития
% ============================================

\chapter{ОБСУЖДЕНИЕ И НАПРАВЛЕНИЯ РАЗВИТИЯ}

\section{Ограничения текущей реализации}

\subsection{Ограничения генеративных возможностей}

Текущая реализация платформы \texttt{sdCloud} обладает рядом ограничений, определяющих границы применимости системы и направления дальнейшего совершенствования.

\textbf{Ограничения больших языковых моделей.} Качество генерируемых моделей напрямую зависит от возможностей используемой LLM. В текущей версии платформы используются в основном бесплатные модели семейства Gemma (30B параметров) на OpenRouter, что накладывает дополнительные ограничения:

\begin{itemize}
    \item \emph{Галлюцинации и фактические ошибки.} LLM могут генерировать структурно корректные, но содержательно ошибочные причинно-следственные связи, нулевые значения констант, пропускать скобки в формулах или использовать неподдерживаемые функции (например, \texttt{MAX} в верхнем регистре).

    \item \emph{Ограничение контекста и внимания.} Даже модели с формально большим контекстом испытывают трудности при генерации моделей с числом переменных более 10--15. Сложные системы требуют декомпозиции или более мощных моделей.

    \item \emph{Недетерминированность и нестабильность free-tier.} Идентичный запрос может давать различные результаты; бесплатные модели могут быть временно недоступны, иметь низкие rate limits и отказывать в генерации. Это затрудняет воспроизводимость и надёжность системы.

    \item \emph{Ограниченные рассуждения и structured output.} 30B-модели хуже справляются с многошаговыми рассуждениями, критикой собственных ответов и строгим соблюдением JSON-схем, что критично для мультиагентных сценариев.

    \item \emph{Ограниченные знания предметных областей.} LLM обладают широкими, но поверхностными знаниями. Для специализированных предметных областей качество генерации снижается.
\end{itemize}

В текущей версии влияние перечисленных ограничений частично компенсируется детерминированным генератором Python-кода, постпроцессингом уравнений, fallback-симуляцией и валидацией графов.

\textbf{Ограничения конвейера генерации и визуализации.} Многоэтапный конвейер, обеспечивающий прозрачность процесса, одновременно вносит дополнительные точки отказа:

\begin{itemize}
    \item Ошибки на ранних этапах (переменные, CLD) распространяются на последующие этапы (SFD, уравнения), поскольку каждый этап использует результат предыдущего.
    \item Межэтапная валидация обнаруживает ошибки, но механизм автоматической коррекции (iterative refinement) требует дополнительных итераций взаимодействия с LLM, что увеличивает время генерации и стоимость.
    \item Конвейер не поддерживает параллельную генерацию альтернативных вариантов модели для одного описания.
    \item Визуализация SFD и CLD на базе React Flow использует упрощённую ручную/круговую раскладку; при большом числе узлов рёбра могут пересекаться, а текст узлов остаётся мелким.
    \item Отсутствует автоматическая калибровка параметров по историческим данным и полноценный анализ чувствительности.
\end{itemize}

\subsection{Ограничения программной реализации}

\textbf{Масштабируемость.} Текущая архитектура предполагает развёртывание на одном сервере. При увеличении нагрузки необходимо внедрение горизонтального масштабирования (кластеризация воркеров Taskiq, репликация PostgreSQL, шардирование Redis).

\textbf{Аутентификация и авторизация.} В текущей реализации используется упрощённая система аутентификации на базе JWT-токенов. Для промышленного использования требуется интеграция с внешними провайдерами идентификации (OAuth 2.0, SAML), системой ролей и разграничением доступа к сессиям и моделям.

\textbf{Обработка ошибок.} Механизм обработки ошибок генерации требует доработки: необходимо реализовать более информативные сообщения об ошибках, систему логирования с разделением по уровням и мониторинг состояния системы.

\textbf{Совместимость форматов.} Текущая реализация поддерживает экспорт моделей в Python-код (основной исполняемый формат) и в текстовый формат Vensim .mdl. Для полноценной интеграции с экосистемой инструментов требуется поддержка формата XMILE, корректная работа с PySD 3.x и возможность импорта/экспота Stella.


\section{Планируемые улучшения}

\subsection{Краткосрочные улучшения текущей версии}

В рамках текущей версии платформы (до перехода к мультиагентной архитектуре) планируется реализация следующих улучшений:

\begin{itemize}
    \item \textbf{Улучшение визуализации графов.} Интеграция библиотек автоматической раскладки (\texttt{dagre}, \texttt{elk.js}, \texttt{cola}) для уменьшения пересечений рёбер в CLD/SFD и улучшения читаемости узлов.
    \item \textbf{Редактор моделей.} Возможность ручного редактирования переменных, уравнений и графов в интерфейсе с немедленным пересчётом Python-кода и симуляции.
    \item \textbf{Полноценная интеграция с PySD 3.x.} Адаптация генератора Python-кода под актуальный API PySD либо полный переход на SciPy-раннер с поддержкой табличных функций, задержек и lookup-функций.
    \item \textbf{Анализ чувствительности.} Автоматический запуск симуляций с вариацией ключевых параметров и визуализация диапазонов поведения модели.
    \item \textbf{Валидация уравнений.} Символьная проверка размерностей, поиск циклов в уравнениях, проверка корректности функций \texttt{integ}.
\end{itemize}

\subsection{Интеграция локальных и платных языковых моделей}

Приоритетным направлением развития является расширение поддержки моделей:

\begin{itemize}
    \item \textbf{Платные модели} (Claude 3.5 Sonnet, GPT-4o-mini, Gemini Flash) для повышения качества сложной генерации, критики и рассуждений.
    \item \textbf{Локальные LLM} через \textbf{Ollama} или \textbf{vLLM} для обеспечения конфиденциальности, offline-режима и снижения операционных затрат.
\end{itemize}

Использование free-tier 30B-моделей в текущей версии является компромиссом для прототипирования; для серьёзной мультиагентной системы и сложных моделей требуются более надёжные модели.

\subsection{Отложенное направление: мультиагентная система}

Разработка полноценной мультиагентной системы (агенты извлечения, построения CLD, генерации уравнений, критика, симуляции и др.) отложена до момента обеспечения:

\begin{itemize}
    \item Надёжного structured output и function calling используемой LLM.
    \item Достаточного контекстного окна для передачи состояния между агентами.
    \item Стабильной работы API без частых rate limits и отказов.
\end{itemize}

В текущей версии функции потенциальных агентов частично реализованы в виде отдельных модулей (\texttt{CLDBuilder}, \texttt{SFDBuilder}, \texttt{GraphValidator}, \texttt{python\_generator.py}), что создаёт основу для будущего перехода к мультиагентной архитектуре.

\subsection{Коллаборативное моделирование и интеграция с данными}

Перспективными направлениями являются:

\begin{itemize}
    \item Совместное редактирование моделей несколькими пользователями в реальном времени.
    \item Версионирование моделей и система комментариев.
    \item Загрузка исторических данных для калибровки параметров и валидации прогнозов.
    \item Публикация моделей в открытый репозиторий.
\end{itemize}


\section{Потенциальные приложения}

\subsection{Образование}

Платформа \texttt{sdCloud} имеет значительный потенциал в образовательной сфере:

\begin{itemize}
    \item \textbf{Обучение системной динамике.} Студенты могут быстро создавать модели на основе текстовых описаний, фокусируясь на анализе структуры и поведения системы, а не на технических аспектах моделирования.
    \item \textbf{Интерактивные лекции.} Преподаватели могут демонстрировать динамику различных систем в реальном времени, просто описывая их словами.
    \item \textbf{Проверка понимания.} Система может генерировать модели по описаниям студентов, позволяя преподавателю оценить глубину понимания предмета.
    \item \textbf{Междисциплинарное обучение.} Специалисты из различных областей (экология, экономика, медицина) могут создавать модели без глубокой подготовки в системной динамике.
\end{itemize}

\subsection{Бизнес-аналитика}

В сфере бизнеса платформа может применяться для:

\begin{itemize}
    \item \textbf{Быстрого прототипирования моделей.} Бизнес-аналитики могут оперативно создавать модели бизнес-процессов для оценки последствий решений.
    \item \textbf{Стратегического планирования.} Моделирование сценариев развития бизнеса с учётом обратных связей и запаздываний.
    \item \textbf{Управление цепями поставок.} Анализ динамики запасов, производства и спроса с выявлением эффектов хлыстовой верёвки (bullwhip effect).
\end{itemize}

\subsection{Анализ государственной политики}

Системная динамика традиционно применяется для анализа последствий государственных решений. Платформа \texttt{sdCloud} может снизить барьер для использования методологии в государственных учреждениях:

\begin{itemize}
    \item Моделирование динамики населения и демографических процессов.
    \item Анализ эпидемиологических сценариев и оценка эффективности вмешательств.
    \item Оценка последствий экологической политики и перехода к устойчивому развитию.
    \item Моделирование динамики транспортных систем и планирование инфраструктуры.
\end{itemize}

\subsection{Научные исследования}

В академической среде платформа может служить инструментом:

\begin{itemize}
    \item Быстрой проверки гипотез о структуре сложных систем.
    \item Автоматизации рутинных этапов моделирования для фокусировки на содержательном анализе.
    \item Создания репозиториев моделей с возможностью воспроизведения и верификации.
    \item Интеграции моделей системной динамики с методами машинного обучения (гибридное моделирование).
\end{itemize}

\subsection{Разработка программного обеспечения}

Методологии, отработанные в проекте \texttt{sdCloud}, имеют потенциал для применения в смежных областях:

\begin{itemize}
    \item Генерация формальных спецификаций из естественно-языковых требований.
    \item Автоматическое построение UML-диаграмм из текстовых описаний систем.
    \item Создание цифровых двойников на основе текстовой документации.
\end{itemize}

Таким образом, результаты настоящего исследования имеют междисциплинарную значимость и могут служить основой для дальнейших разработок на стыке искусственного интеллекта, системной динамики и программной инженерии.
